\documentclass[11pt]{IEEEtran}
\usepackage[left=1.5in, right=1in, top=1in, bottom=1in]{geometry}
\usepackage{setspace}
\usepackage{graphicx}
\doublespacing

%%%%%%%%%%%%%%%%%%%%%%%%%%%%%%%%%%%%
%%%%%%%%%%%% DOC START %%%%%%%%%%%%%
%%%%%%%%%%%%%%%%%%%%%%%%%%%%%%%%%%%%
\begin{document}

%%%%%%%%%%%%%%%%%%%%%%%%%%%%%%%%%%%%
%%%%%%% REMOVE CHECK SECTION %%%%%%%
%%%%%%%%%%%%%%%%%%%%%%%%%%%%%%%%%%%%

\onecolumn
\singlespacing
\vspace{-5mm}
\hrulefill \textbf{Remove this section} \hrulefill
\par
As per syllabus, we need to ensure the following: 
\par
\begin{itemize}

\item Don't Plagiarize (Only one of us will turn in the paper electronically). Perhaps we will both need printed form submissions? *Tables and Figures not required for electronic submission

\item 15 Pages of text which does not include:
\begin{itemize}
    \item Cover Page, Table of Contents, List of Tables, List of figures (intro stuff)
    \item Conclusion
    \item References
    \item Bibliography
    \item Tables.  Possibly tables and figures as the syllabus may be in error here.
\end{itemize}
\item 11 Point font
\item Double spaced, with Triple spacing after headings
\item Margins: 1.5" left, 1" top right bottom (binding standard)
\item Format:
\begin{itemize}
\item Cover Page: title, abstract, course name and number, date
\item Table of Contents: page numbers for major headings
\item List of Tables: page numbers w/ captions
\item List of Figures: Page numbers w/ captions
\item Intro
\item Body w/ figures and tables
\item conclusion
\item references
\item bibliography
\item tables presented in order together
\item figures presented in order together
\end{itemize}
\item "Typical Errors"
\begin{itemize}
\item Did not reference sources in text/tables/figures
\item margins are wrong, spacing and triple spacing
\item grammar / spelling
\item did not follow format
\end{itemize}
\end{itemize}
\doublespacing

\begin{titlepage}
    \begin{center}
        \vspace*{-20mm}
        \huge
        \textbf{Data Flow Architecture:\\ \vspace{-3mm} Benefits, Limitations and Future Applications}\\
         \large
        
        \vspace{0.3 cm}
        \textbf{
        Jake Baker (jacobwilliambaker@csus.edu) \\
        \vspace{-2mm}
        Jesse Harkness (jesseharkness@csus.edu)} \\ 
        \vspace{-2mm}
        EEE285 - Micro-Computer System Design I \\
        \vspace{-2mm}
        INSERT/ACTUAL/DATE\\
        \vspace{10mm}
        \includegraphics[width=0.3\textwidth]{Picture Files/1200px-California_State_University,_Sacramento_seal.svg.png}
        \vfill
        \textbf{Abstract}\\
        \raggedright\setlength{\parindent}{3em}
        The automation of everyday tasks is the magnum opus of computer and information technologies; as computer processing ability continues to increase, the next ubiquitous technological advance will be through artificial intelligence (AI).  But the level of computation required for AI applications comes at a cost; increasing transistor density is no longer a feasible method for increasing compute power.  In contrast to the traditional CPU architecture which relies on smaller and denser spreads of transistors, Data Flow Architecture (DFA) is a hardware design which maximizes compute potential by reducing data movement, and increasing work efficiency through highly specialized processors.  By exploiting the availability of data, rather than the sequence of instructions, DFA improves concurrency of operations and has the potential to maximize throughput.  Comparing the similarities and differences of the dataflow paradigm with a traditional CPU highlights the benefits of DFA, some of the inherent challenges, as well as the future possibilties for this hardware architecture in the AI space.     
\par


    \end{center}
\end{titlepage}
\pagenumbering{roman}
\clearpage
\pagenumbering{arabic}

\renewcommand\IEEEkeywordsname{Keyword Index}
\begin{IEEEkeywords}
Artificial Intelligence, compute, control flow architecture,  CPU, data flow design, deep learning, dynamic dataflow, GPU, machine learning, machine vision, neural networks, scheduling, tensor flow 
\end{IEEEkeywords}

\clearpage
\section{Section 1}
Companies like SambaNova are at the forefront of this hardware technology, creating chips which can support AI models with up to 5 trillion parameters [4]; however, in order for this impressive technology to become trending, manufacturers need to overcome a host of challenges blocking the next generation of AI fueled technology.  While the aforementioned processors at SambaNova are undoubtedly powerful, the lack of software support for most data flow architectures is creating a void between what's possible, versus what's practical.  Data Flow Architectures have the potential to provide immense compute capability for the field of AI and Machine Learning; however the fledgling software ecosystem to support this architecture creates a unique challenge for its practical use, thus stifling its potential mass consumerization.  More traditional computer architectures, like the Von Neumann CPU architecture and GPU architecture, have a rich software backbone making it a more pragmatic choice for many companies with products the AI/ML space.  Overcoming the current nascent software landscape will be a necessary step in the future of data flow architectures for AI/ML applications. 

\section{Section 2}
Companies like SambaNova are at the forefront of this hardware technology, creating chips which can support AI models with up to 5 trillion parameters [4]; however, in order for this impressive technology to become trending, manufacturers need to overcome a host of challenges blocking the next generation of AI fueled technology.  While the aforementioned processors at SambaNova are undoubtedly powerful, the lack of software support for most data flow architectures is creating a void between what's possible, versus what's practical.  Data Flow Architectures have the potential to provide immense compute capability for the field of AI and Machine Learning; however the fledgling software ecosystem to support this architecture creates a unique challenge for its practical use, thus stifling its potential mass consumerization.  More traditional computer architectures, like the Von Neumann CPU architecture and GPU architecture, have a rich software backbone making it a more pragmatic choice for many companies with products the AI/ML space.  Overcoming the current nascent software landscape will be a necessary step in the future of data flow architectures for AI/ML applications. 

\begin{thebibliography}{00}

\bibitem{b1} SoftwareMill, "Technology Trends," SoftwareMill Blog. [Online]. Available: https://softwaremill.com/technology-trends/. [Accessed: 9-26-2023].
.

\bibitem{b2} "Microservices usage worldwide 2020, by organization size," Statista, 2020. [Online]. Available: https://www.statista.com/statistics/1236823/microservices-usage-per-organization-size/. [Accessed: 9-27-2023].


\bibitem{b3} A. Bitar - Groq, "Groq’s Software-Defined Hardware for Dataflow Compute," YouTube, 27 Feb. 2021. [Video]. Available: https://www.youtube.com/watch?v=nP\_9lkDsD-E\&t=919s. [Accessed: 9-26-2023].

\bibitem{b4} E. Gent, "SambaNova’s New Chip Means GPTs for Everyone," IEEE Spectrum, [Online]. Available: https://spectrum.ieee.org/ai-chip-sambanova. [Accessed: 09-25-2023].

\end{thebibliography}


\end{document}
